\section{Introduction}
\label{sec:intro}

Text-based Person Search (TBPS) is a cross-modal retrieval task that locates target individuals in large image galleries using natural language descriptions~\cite{1}. As a bridge between Person Re-identification (Re-ID) and text-to-image retrieval, TBPS allows for more flexible and descriptive queries than traditional image-based methods. However, TBPS remains challenging due to redundant visual features~\cite{2,3} and the inherent noise in cross-modal correspondence~\cite{xie2025dynamic}. Furthermore, the lack of adaptive mechanisms for handling ambiguous queries during inference limits practical application.

\begin{figure}[h!]
\centering
\includegraphics[width=\linewidth]{fig1.pdf}
\caption{A comparison of TBPS.(a) \textbf{Existing Method:} The search is performed directly using the original text query.(b) \textbf{Our IQE Approach:} We improve the query at inference time without retraining. First, our method finds a relevant anchor image. Then, an MLLM learns key visual details from this image through a Q\&A process. Finally, it fuses these details with the original text to create an improved query and re-ranks the search results.}
\label{fig:iqe}
\end{figure}

Most existing methods improve cross-modal alignment using fine-grained annotations or large pre-trained models like CLIP~\cite{4}. While integrated planning and scheduling frameworks have proven effective in optimizing complex systems~\cite{liu2025integratedplanningmachinelevelscheduling}, TBPS methods still mainly rely on \textit{passive alignment}—matching features in the embedding space without actively improving the inputs. Some approaches explore hierarchical visual perception to achieve deeper alignment~\cite{xie2026delvingdeeperhierarchicalvisual,xie2026hvdhumanvisiondrivenvideo}, yet they also suffer from this passive limitation. This leads to two key problems: (1) Methods such as CFAM~\cite{5} depend on manual annotations, but their performance drops when annotations are scarce; (2) Methods like DM-Adapter~\cite{6} and MGRL~\cite{7} cannot refine brief or ambiguous queries at inference, resulting in poor results for incomplete descriptions.

Some recent works have attempted active refinement or interaction. For instance, chat-driven mechanisms have been proposed to generate text for retrieval~\cite{xie2025chat}, though purely generative approaches may hallucinate attributes. Other works like TVFR~\cite{8} are sensitive to the initial query, and AUL~\cite{9} mainly helps during training. \textit{These limitations motivate us to design a method that achieves fine-grained cross-modal alignment during training and actively refines ambiguous or incomplete queries at inference.}

To solve the aforementioned problems, we propose \textbf{CONQUER}, a two-stage framework as illustrated in Fig.~\ref{fig:iqe}. During training, our \textbf{Context-Aware Representation Enhancement (CARE)} module learns robust cross-modal embeddings. Subsequently, at the inference stage, our \textbf{Interactive Query Enhancement (IQE)} module acts as a plug-and-play enhancement that adaptively refines user queries. By leveraging fine-grained attributes extracted from high-confidence candidate images, IQE enriches ambiguous or underspecified descriptions. The enhanced information is then fused with the original query, enabling more accurate re-ranking of the gallery and improved retrieval results.

Our main contributions are as follows:
\begin{itemize}
    \item We introduce CARE, a context-aware representation enhancement module that addresses the cross-modal gap by combining multi-granularity representation encoding, complementary pair mining, and context-guided optimal matching via Optimal Transport.
    \item We develop IQE, an interactive query enhancement strategy that adaptively refines ambiguous or incomplete user queries during inference by leveraging attribute information from candidate images.
    
    \item Extensive experiments show that CONQUER achieves state-of-the-art performance on several TBPS benchmarks and demonstrates strong robustness in cross-domain and incomplete-query scenarios.
\end{itemize}